% Figure: Conservation Laws Overview
% Chapter 3: Constitutive Equations and Conservation Laws
\begin{figure}[htbp]
\centering
\begin{tikzpicture}[
    domain box/.style={draw=UofSCBlack, fill=#1!15, rectangle,
                       minimum height=2.2cm, minimum width=4.5cm,
                       sharp corners, font=\small, align=center},
    law label/.style={font=\bfseries\small, #1},
    equation/.style={font=\footnotesize, UofSC70Black}
]
    % Central concept
    \node[draw=UofSCGarnet, fill=UofSCGarnet!10, rectangle,
          minimum height=1.5cm, minimum width=5cm,
          line width=1.5pt, sharp corners, font=\bfseries]
          (center) at (0, 0) {Conservation Laws};

    % Mechanical domain (top left)
    \node[domain box=UofSCAtlantic, anchor=south east] (mech) at (-0.5, 1.2) {
        \textbf{Mechanical}\\[3pt]
        Newton's 2nd Law\\[2pt]
        $\sum F = m\ddot{x}$\\
        $\sum \tau = J\ddot{\theta}$
    };
    \node[law label=UofSCAtlantic, above=0.1cm of mech] {Momentum Conservation};

    % Electrical domain (top right)
    \node[domain box=UofSCHorseshoe, anchor=south west] (elec) at (0.5, 1.2) {
        \textbf{Electrical}\\[3pt]
        Kirchhoff's Laws\\[2pt]
        $\sum i = 0$ (KCL)\\
        $\sum v = 0$ (KVL)
    };
    \node[law label=UofSCHorseshoe, above=0.1cm of elec] {Charge Conservation};

    % Thermal domain (bottom left)
    \node[domain box=UofSCRose, anchor=north east] (therm) at (-0.5, -1.2) {
        \textbf{Thermal}\\[3pt]
        Energy Balance\\[2pt]
        $C\dfrac{dT}{dt} = \dot{Q}_{\text{in}} - \dot{Q}_{\text{out}}$
    };
    \node[law label=UofSCRose, below=0.1cm of therm] {Energy Conservation};

    % Fluid domain (bottom right)
    \node[domain box=UofSCCongaree, anchor=north west] (fluid) at (0.5, -1.2) {
        \textbf{Fluid}\\[3pt]
        Continuity Equation\\[2pt]
        $A_1 v_1 = A_2 v_2$\\
        $\dfrac{d(\rho V)}{dt} = \dot{m}_{\text{in}} - \dot{m}_{\text{out}}$
    };
    \node[law label=UofSCCongaree, below=0.1cm of fluid] {Mass Conservation};

    % Connecting lines
    \draw[thick, UofSCAtlantic] (center.north west) -- (mech.south east);
    \draw[thick, UofSCHorseshoe] (center.north east) -- (elec.south west);
    \draw[thick, UofSCRose] (center.south west) -- (therm.north east);
    \draw[thick, UofSCCongaree] (center.south east) -- (fluid.north west);

    % Analogy connections (dashed, showing cross-domain similarities)
    \draw[dashed, UofSC50Black, <->] (mech.east) -- (elec.west)
        node[midway, above, font=\tiny] {Analogy};
    \draw[dashed, UofSC50Black, <->] (therm.east) -- (fluid.west)
        node[midway, above, font=\tiny] {Analogy};

\end{tikzpicture}
\caption{The four fundamental conservation laws and their application domains. Each physical domain has corresponding conservation principles that govern system dynamics. Cross-domain analogies enable unified analysis techniques.}
\label{fig:conservation_overview}
\end{figure}
